% ==============================================================================
% SLIDES - SEMANA 16: CONSOLIDAÇÃO FINAL DO PDGTI, CHECKLIST & GOVERNANÇA CONTÍNUA
% ==============================================================================

% 1. CAPA
\senaicover
  {Consolidação Final do PDGTI}
  {Checklist de Qualidade, Governança Contínua (PDCA) e Entrega Final}
  {Unidade Curricular: Governança de TI}
  {Prof. André Souza}

% 2. AGENDA
\begin{senaiframe}{Visão Geral \& Agenda}{Estrutura da Aula}
  \begin{columns}[t]
    \begin{column}{0.48\textwidth}
      \begin{tcolorbox}[senaibluecard, title={1. Visão Sistêmica}]
        \begin{itemize}
          \item Coerência inter-capítulos (Cap. 1 ao 9).
          \item Encadeamento de diagnóstico a roadmap.
        \end{itemize}
      \end{tcolorbox}
      \vspace{4pt}
      \begin{tcolorbox}[senaiambercard, title={3. Governança Contínua}]
        \begin{itemize}
          \item Ciclo PDCA no PDGTI.
          \item Ritos e revisões anuais do Comitê.
        \end{itemize}
      \end{tcolorbox}
    \end{column}
    \begin{column}{0.48\textwidth}
      \begin{tcolorbox}[senairedcard, title={2. Checklist de Qualidade}]
        \begin{itemize}
          \item Verificação de rigor técnico e formatação.
          \item Matriz RACI, SWOT, BSC e Riscos.
        \end{itemize}
      \end{tcolorbox}
      \vspace{4pt}
      \begin{tcolorbox}[senaigreencard, title={4. Entrega Final}]
        \begin{itemize}
          \item Submissão do PDGTI Completo (30%).
          \item Encerramento dos entregáveis.
        \end{itemize}
      \end{tcolorbox}
    \end{column}
  \end{columns}
\end{senaiframe}

% 3. ARQUITETURA ORGANIZACIONAL
\begin{senaiframe}{Teoria • Módulo 1}{A Arquitetura Organizacional do PDGTI}
  \begin{columns}[t]
    \begin{column}{0.48\textwidth}
      \begin{tcolorbox}[senaibluecard, title={Artefato Vivo}]
        \begin{itemize}
          \item O PDGTI é o plano mestre de transformação digital da empresa.
          \item Não é um relatório isolado, mas uma diretriz contínua.
        \end{itemize}
      \end{tcolorbox}
    \end{column}
    \begin{column}{0.48\textwidth}
      \begin{tcolorbox}[senaicard, title={Harmonia dos Conteúdos}]
        \begin{itemize}
          \item Cada capítulo sustenta e justifica o seguinte.
        \end{itemize}
      \end{tcolorbox}
    \end{column}
  \end{columns}
\end{senaiframe}

% 4. VISÃO SISTÊMICA
\begin{senaiframe}{Teoria • Módulo 1}{Visão Sistêmica: O Encadeamento Lógico}
  \begin{columns}[t]
    \begin{column}{0.48\textwidth}
      \begin{tcolorbox}[senaibluecard, title={Base \& Diagnóstico}]
        \begin{itemize}
          \item Caracterização (Cap 1) -> Diagnóstico SWOT (Cap 2) -> Estratégia BSC (Cap 3).
        \end{itemize}
      \end{tcolorbox}
    \end{column}
    \begin{column}{0.48\textwidth}
      \begin{tcolorbox}[senaigreencard, title={Controle \& Execução}]
        \begin{itemize}
          \item RACI (Cap 4) -> KPIs (Cap 5) -> Riscos (Cap 6) -> Frameworks (Cap 7-8) -> Roadmap (Cap 9).
        \end{itemize}
      \end{tcolorbox}
    \end{column}
  \end{columns}
\end{senaiframe}

% 5. CAPÍTULOS 1 E 2
\begin{senaiframe}{Teoria • Módulo 1}{Do Contexto ao Diagnóstico (Capítulos 1 \& 2)}
  \begin{columns}[t]
    \begin{column}{0.48\textwidth}
      \begin{tcolorbox}[senaicard, title={Capítulo 1: Onde Estamos}]
        \begin{itemize}
          \item Porte, setor, missão e estrutura atual de TI.
        \end{itemize}
      \end{tcolorbox}
    \end{column}
    \begin{column}{0.48\textwidth}
      \begin{tcolorbox}[senairedcard, title={Capítulo 2: O que nos Dói}]
        \begin{itemize}
          \item SWOT de TI, gargalos operacionais e nível de maturidade.
        \end{itemize}
      \end{tcolorbox}
    \end{column}
  \end{columns}
\end{senaiframe}

% 6. CAPÍTULOS 3 E 4
\begin{senaiframe}{Teoria • Módulo 1}{Da Estratégia à Decisão (Capítulos 3 \& 4)}
  \begin{columns}[t]
    \begin{column}{0.48\textwidth}
      \begin{tcolorbox}[senaibluecard, title={Capítulo 3: Aonde Queremos Ir}]
        \begin{itemize}
          \item Mapa Estratégico do IT BSC e Portfólio de TI.
        \end{itemize}
      \end{tcolorbox}
    \end{column}
    \begin{column}{0.48\textwidth}
      \begin{tcolorbox}[senaiambercard, title={Capítulo 4: Quem Decide}]
        \begin{itemize}
          \item Comitê Estratégico, Matriz RACI e Políticas de TI.
        \end{itemize}
      \end{tcolorbox}
    \end{column}
  \end{columns}
\end{senaiframe}

% 7. CAPÍTULOS 5 E 6
\begin{senaiframe}{Teoria • Módulo 1}{Da Medição à Proteção (Capítulos 5 \& 6)}
  \begin{columns}[t]
    \begin{column}{0.48\textwidth}
      \begin{tcolorbox}[senaigreencard, title={Capítulo 5: Como Medimos}]
        \begin{itemize}
          \item KPIs do Negócio e TI, Baseline e Dashboard.
        \end{itemize}
      \end{tcolorbox}
    \end{column}
    \begin{column}{0.48\textwidth}
      \begin{tcolorbox}[senairedcard, title={Capítulo 6: Como nos Protegemos}]
        \begin{itemize}
          \item Matriz de Riscos 5x5, Mitigação e LGPD.
        \end{itemize}
      \end{tcolorbox}
    \end{column}
  \end{columns}
\end{senaiframe}

% 8. CAPÍTULOS 7 E 8
\begin{senaiframe}{Teoria • Módulo 1}{Dos Frameworks à Qualidade (Capítulos 7 \& 8)}
  \begin{columns}[t]
    \begin{column}{0.48\textwidth}
      \begin{tcolorbox}[senaibluecard, title={Capítulo 7: Melhores Práticas}]
        \begin{itemize}
          \item Seleção de COBIT 2019, ITIL v4 e ISO 38500.
        \end{itemize}
      \end{tcolorbox}
    \end{column}
    \begin{column}{0.48\textwidth}
      \begin{tcolorbox}[senaiambercard, title={Capítulo 8: Qualidade de Dev}]
        \begin{itemize}
          \item Maturidade CMMI/MPS-SW e ISO 25010.
        \end{itemize}
      \end{tcolorbox}
    \end{column}
  \end{columns}
\end{senaiframe}

% 9. CAPÍTULO 9
\begin{senaiframe}{Teoria • Módulo 1}{Da Execução ao Valor (Capítulo 9)}
  \begin{columns}[t]
    \begin{column}{0.48\textwidth}
      \begin{tcolorbox}[senaigreencard, title={Capítulo 9: Quando e Quanto}]
        \begin{itemize}
          \item Roadmap em 3 Ondas, Orçamento CAPEX/OPEX.
          \item Fatores Críticos de Sucesso e Geração de Valor.
        \end{itemize}
      \end{tcolorbox}
    \end{column}
    \begin{column}{0.48\textwidth}
      \begin{tcolorbox}[senaicard, title={Síntese Final}]
        \begin{itemize}
          \item Fechamento do plano de investimentos de 3 anos.
        \end{itemize}
      \end{tcolorbox}
    \end{column}
  \end{columns}
\end{senaiframe}

% 10. GOVERNANÇA CONTÍNUA
\begin{senaiframe}{Teoria • Módulo 2}{O Conceito de Governança Contínua}
  \begin{columns}[t]
    \begin{column}{0.48\textwidth}
      \begin{tcolorbox}[senaibluecard, title={Instrumento Vivo}]
        \begin{itemize}
          \item O PDGTI deve ser revisado anualmente.
          \item Reajuste de metas conforme mudanças de mercado.
        \end{itemize}
      \end{tcolorbox}
    \end{column}
    \begin{column}{0.48\textwidth}
      \begin{tcolorbox}[senaicard, title={Evitar o Engavetamento}]
        \begin{itemize}
          \item Acompanhamento mensal das metas pelo Comitê.
        \end{itemize}
      \end{tcolorbox}
    \end{column}
  \end{columns}
\end{senaiframe}

% 11. O CICLO PDCA NO PDGTI
\begin{senaiframe}{Teoria • Módulo 2}{O Ciclo PDCA Aplicado ao PDGTI}
  \begin{columns}[t]
    \begin{column}{0.48\textwidth}
      \begin{tcolorbox}[senaibluecard, title={Plan \& Do}]
        \begin{itemize}
          \item \textbf{Plan:} Elaborar/Revisar o PDGTI e metas.
          \item \textbf{Do:} Executar as ondas do roadmap e políticas.
        \end{itemize}
      \end{tcolorbox}
    \end{column}
    \begin{column}{0.48\textwidth}
      \begin{tcolorbox}[senaigreencard, title={Check \& Act}]
        \begin{itemize}
          \item \textbf{Check:} Acompanhar dashboards e auditorias.
          \item \textbf{Act:} Corrigir desvios no Comitê de TI.
        \end{itemize}
      \end{tcolorbox}
    \end{column}
  \end{columns}
\end{senaiframe}

% 12. RITOS DO COMITÊ DE TI
\begin{senaiframe}{Teoria • Módulo 2}{Ritos do Comitê Estratégico de TI}
  \begin{columns}[t]
    \begin{column}{0.48\textwidth}
      \begin{tcolorbox}[senaibluecard, title={Pauta Mensal}]
        \begin{itemize}
          \item Revisão dos KPIs do dashboard e incidentes graves.
        \end{itemize}
      \end{tcolorbox}
    \end{column}
    \begin{column}{0.48\textwidth}
      \begin{tcolorbox}[senaiambercard, title={Pauta Trimestral/Anual}]
        \begin{itemize}
          \item Avaliação do portfólio de projetos e revisão do PDGTI.
        \end{itemize}
      \end{tcolorbox}
    \end{column}
  \end{columns}
\end{senaiframe}

% 13. ESTUDO DE CASO PRÁTICO
\begin{senaiframe}{Estudo de Caso}{Auditoria de Qualidade em Artefato de PDGTI}
  \begin{columns}[t]
    \begin{column}{0.48\textwidth}
      \begin{tcolorbox}[senairedcard, title={Falha de Auditoria}]
        \begin{itemize}
          \item Matriz RACI com 2 'A's em alguns processos.
          \item Incompatibilidade corrigida na revisão final.
        \end{itemize}
      \end{tcolorbox}
    \end{column}
    \begin{column}{0.48\textwidth}
      \begin{tcolorbox}[senaigreencard, title={Aprovação Executiva}]
        \begin{itemize}
          \item Documento 100\% auditável e aprovado pelo Conselho.
        \end{itemize}
      \end{tcolorbox}
    \end{column}
  \end{columns}
\end{senaiframe}

% 14. REFERÊNCIAS BIBLIOGRÁFICAS
\begin{senaiframe}{Referências Bibliográficas}{Fontes \& Leitura Recomendada}
  \begin{tcolorbox}[senaicard, title={Obras de Referência}]
    \begin{itemize}
      \item \textbf{FERNANDES, Aguinaldo Aragon; ABREU, Vladimir Ferraz.} *Implantando a governança de TI: da estratégia à gestão de processos e serviços.* 4. ed. Rio de Janeiro: Brasport, 2014.
      \item \textbf{WEILL, Peter; ROSS, Jeanne W.} *Governança de Tecnologia da Informação.* Rio de Janeiro: Elsevier, 2004.
    \end{itemize}
  \end{tcolorbox}
\end{senaiframe}

% 15. APLICAÇÃO NO PDGTI
\begin{senaiframe}{Aplicação no PDGTI}{Instruções para a ENTREGA FINAL (27/Nov)}
  \vspace{0.2cm}
  \begin{tcolorbox}[senaigreencard, title={Entrega Final do PDGTI Completo (Peso: 30%)}]
    \begin{itemize}
      \item Submissão do documento consolidado com os **Capítulos 1 a 9**.
      \item Incorporar todas as correções e devolutivas das Entregas Parciais 1, 2 e 3.
      \item Data limite impreterível: **27 de Novembro**.
    \end{itemize}
  \end{tcolorbox}
\end{senaiframe}
